\chapter{Per-scene min-max, and a library with nowhere to stand}\label{ch:hslnorm}
\def\tbrk{\hfil\penalty50\hfilneg}% local: a line may break before a provenance tag

\begin{onething}
Per-scene min-max scaling makes a pixel's input depend on the brightest and darkest pixels of its scene, so the same
material reads differently in every scene, one bright anomaly changes every other pixel, and a library spectrum, which
has no scene, has no defined input at all; measured on four PRISMA scenes, HyperSL's embeddings did not name man-made
targets better than plain spectral angle, and one vegetation spectrum absorbed 60\% of a scene's anomalies.
\end{onething}

Library matching (\cref{sec:material:naming}, \cref{ch:namer}) asks whether a pixel looks like a laboratory spectrum.
That question has a precondition so basic it is easy to miss: \emph{the pixel and the library spectrum must be put on
the same footing before they are compared}. This chapter shows, on numbers small enough to do by hand, that HyperSL's
input scaling cannot meet that precondition, and then shows what happened when the team tried anyway.

\begin{analogy}
Marks graded on a curve. In each class the weakest student is given 0 and the strongest 100, and everyone else is
placed proportionally between. A mark of 75 means one thing in a class of strong students and another in a weak one;
a single brilliant newcomer pushes everyone else's mark down; and an external candidate who sat the exam alone has no
curved mark at all, because there is no class to curve against. Comparing that candidate with a student in some class
means inventing a class for them. The library spectrum is the external candidate.
\end{analogy}

%======================================================================
\section{The scaling, as a formula}\label{sec:hslnorm:formula}

\begin{definition}[Per-scene, per-band min-max scaling]
For a scene with pixels $\vect{x}^{(1)},\dots,\vect{x}^{(N)}$ and band $b$, let
$m_b=\min_n x^{(n)}_b$ and $M_b=\max_n x^{(n)}_b$. Each pixel's value in band $b$ becomes
\begin{equation}
x'_b=\frac{x_b-m_b}{M_b-m_b}. \label{eq:hslnorm:minmax}
\end{equation}
This is \cref{eq:hslwhat:minmax}, the operation \code{sklearn.preprocessing.minmax_scale} applies to each column of the
pixels-by-bands table in \file{upstream/mutihsi.py}.
\end{definition}

Rewrite it to see its shape:
\begin{align}
x'_b &= \frac{x_b-m_b}{M_b-m_b} \notag\\
&= \frac{1}{M_b-m_b}\,x_b-\frac{m_b}{M_b-m_b} && \text{split the fraction}\notag\\
&= s_b\,x_b+c_b, \quad s_b=\frac{1}{M_b-m_b},\ c_b=-\frac{m_b}{M_b-m_b}. \label{eq:hslnorm:affine}
\end{align}
\begin{definition}[Per-band affine map]
A map that multiplies band $b$ by a number $s_b$ and adds a number $c_b$, with each band allowed its own $s_b$ and
$c_b$, is a \emph{per-band affine map}\index{per-band affine map}.
\end{definition}
So min-max is a per-band affine map, and so is the fixed z-score our own models use
(\code{WavelengthStats}, \cref{sec:specdecoder:traps}): $z_b=(x_b-\mu(\lambda_b))/\sigma(\lambda_b)$, which is
\cref{eq:hslnorm:affine} with $s_b=1/\sigma(\lambda_b)$ and $c_b=-\mu(\lambda_b)/\sigma(\lambda_b)$.

\begin{keypoint}
The two maps have the same form. The difference is where $s_b$ and $c_b$ come from: from the \emph{scene} for min-max,
from a \emph{fixed curve} stored in the checkpoint for \code{WavelengthStats}. Everything in this chapter follows from
that one difference.
\end{keypoint}

%======================================================================
\section{A three-pixel scene, by hand}\label{sec:hslnorm:toy}

Two bands are enough. Scene $S_1$ has three pixels:
\[
\vect{p}_1=(0.10,\ 0.30),\qquad \vect{p}_2=(0.20,\ 0.50),\qquad \vect{p}_3=(0.30,\ 0.40)\quad\tagA .
\]

\begin{worked}[label={we:hslnorm:s1}]{Scaling scene $S_1$}
Band 1 holds 0.10, 0.20, 0.30, so $m_1=0.10$, $M_1=0.30$, $M_1-m_1=0.20$. Band 2 holds 0.30, 0.50, 0.40, so $m_2=0.30$,
$M_2=0.50$, $M_2-m_2=0.20$.
\begin{align*}
\vect{p}'_1 &= \left(\frac{0.10-0.10}{0.20},\ \frac{0.30-0.30}{0.20}\right) = (0,\ 0),\\
\vect{p}'_2 &= \left(\frac{0.20-0.10}{0.20},\ \frac{0.50-0.30}{0.20}\right) = \left(\frac{0.10}{0.20},\ \frac{0.20}{0.20}\right) = (0.5,\ 1.0),\\
\vect{p}'_3 &= \left(\frac{0.30-0.10}{0.20},\ \frac{0.40-0.30}{0.20}\right) = \left(\frac{0.20}{0.20},\ \frac{0.10}{0.20}\right) = (1.0,\ 0.5).
\end{align*}
\end{worked}

\subsection{(a) The same spectrum, two different inputs}

Scene $S_2$ contains the very same spectrum $(0.20,0.50)$, beside two others: $(0.05,0.20)$ and $(0.40,0.60)$\tagA.

\begin{worked}[label={we:hslnorm:s2}]{The same spectrum in scene $S_2$}
Band 1: $m_1=0.05$, $M_1=0.40$, $M_1-m_1=0.35$. Band 2: $m_2=0.20$, $M_2=0.60$, $M_2-m_2=0.40$.
\begin{align*}
x'_1 &= \frac{0.20-0.05}{0.35} = \frac{0.15}{0.35} = 0.4286,\\
x'_2 &= \frac{0.50-0.20}{0.40} = \frac{0.30}{0.40} = 0.75 .
\end{align*}
In $S_1$ the spectrum $(0.20,0.50)$ entered the model as $(0.5,1.0)$; in $S_2$ as $(0.4286,0.75)$\tagC.
\end{worked}
\textbf{One material, two inputs.} A model fed these will produce two different embeddings for the same surface, and
nothing in its training told it the two are the same thing: the scenes' extremes, not the material, decided the
difference.

\subsection{(b) One bright anomaly moves everybody}

Add a bright fourth pixel, $\vect{a}=(0.60,0.90)$, to $S_1$\tagA: a metal roof, say.

\begin{worked}[label={we:hslnorm:anom}]{Scene $S_1$ after a bright pixel arrives}
The minima do not change ($m_1=0.10$, $m_2=0.30$). The maxima become $M_1=0.60$ and $M_2=0.90$, so the spans are
$0.60-0.10=0.50$ and $0.90-0.30=0.60$.
\begin{align*}
\vect{p}'_2 &= \left(\frac{0.20-0.10}{0.50},\ \frac{0.50-0.30}{0.60}\right)=\left(\frac{0.10}{0.50},\ \frac{0.20}{0.60}\right)=(0.2,\ 0.3333),\\
\vect{p}'_3 &= \left(\frac{0.30-0.10}{0.50},\ \frac{0.40-0.30}{0.60}\right)=\left(\frac{0.20}{0.50},\ \frac{0.10}{0.60}\right)=(0.4,\ 0.1667),\\
\vect{a}' &= \left(\frac{0.50}{0.50},\ \frac{0.60}{0.60}\right)=(1,\ 1).
\end{align*}
$\vect{p}_2$ was $(0.5,1.0)$ and is now $(0.2,0.3333)$; $\vect{p}_3$ was $(1.0,0.5)$ and is now $(0.4,0.1667)$\tagC. Not
one of their reflectances changed.
\end{worked}

For detection this is perverse. The arrival of an anomaly rewrites the input of every \emph{ordinary} pixel in the
scene, and the anomaly itself is pinned to the corner $(1,1)$ of the box the model was trained in. A real scene has a
million pixels and the extremes are usually set by a handful of them: glint, a metal roof, a fire, a cloud. Those
handful decide the frame for the million.

\subsection{(c) A library spectrum has no frame}

A laboratory spectrum $\vect{t}=(0.70,0.20)$\tagA{} was measured alone. It has no scene, so \cref{eq:hslnorm:minmax}
has no $m_b$ and $M_b$ to use. There are three ways to force it in, and each fails.

\paragraph{Borrow the scene's frame.} Use $S_1$'s $m$ and $M$:
\begin{align*}
t'_1 &= \frac{0.70-0.10}{0.20} = \frac{0.60}{0.20} = 3.0, & t'_2 &= \frac{0.20-0.30}{0.20} = \frac{-0.10}{0.20} = -0.5 .
\end{align*}
Both lie outside $[0,1]$, the range every training input occupied, and outside the range the sigmoid output can
reach (\cref{tab:hslwhat:trace}). Clipping gives $(1,0)$ and throws away how far outside it was. On real data this is
not rare: the team's first attempt clipped 41\% of all ECOSTRESS library values, and the clipping fell mostly on
natural entries\tbrk\tagMr{\texttt{4bbf91d}: hypersl\_model.py docstring}.

\paragraph{Put the spectrum into the scene.} Append $\vect{t}$ to $S_1$ as a fourth pixel and recompute. Now
$m=(0.10,0.20)$ and $M=(0.70,0.50)$, $\vect{t}'=(1,0)$, and $\vect{p}_1$ moves from $(0,0)$ to $(0,0.3333)$\tagC. The
library entry has changed the scene it is being compared against. With 2,000 library entries there would be 2,000
different frames.

\paragraph{Scale each spectrum by itself.} Use each spectrum's own minimum and maximum over its bands (``per-curve'').
That treats pixel and library alike, but it is not what the model was trained on: the validation measured the loss
roughly doubling under it (row~I in \cref{sec:hslwhat:repro}).

\begin{definition}[Circular]
A procedure is \emph{circular}\index{circular procedure} when its answer is needed to set up its own input. Borrowing the
scene's frame is circular for matching: a scene's extremes in each band are typically set by its unusual pixels, the
very ones being named, so the frame the library is placed in depends on what the target is.
\end{definition}

\subsection{(d) Spectral angle does not survive the map}

\begin{definition}[Spectral angle]
The angle between two spectra $\vect{u}$ and $\vect{v}$ (\cref{ch:rx}, \cref{ch:material}) is
\[
\theta=\arccos\frac{\vect{u}\cdot\vect{v}}{\norm{\vect{u}}\,\norm{\vect{v}}}.
\]
It ignores overall brightness: $\theta(k\vect{u},\vect{v})=\theta(\vect{u},\vect{v})$ for any $k>0$. That is why
library matching uses it.
\end{definition}

A tempting hope: the angle ignores brightness, so perhaps it ignores the min-max map too. It does not.

\begin{worked}[label={we:hslnorm:angle}]{The angle between $\vect{p}_2$ and $\vect{p}_3$, before and after}
\emph{Before}, in reflectance:
\begin{align*}
\vect{p}_2\cdot\vect{p}_3 &= 0.20\times0.30+0.50\times0.40 = 0.06+0.20 = 0.26,\\
\norm{\vect{p}_2} &= \sqrt{0.20^2+0.50^2} = \sqrt{0.04+0.25} = \sqrt{0.29} = 0.5385,\\
\norm{\vect{p}_3} &= \sqrt{0.30^2+0.40^2} = \sqrt{0.09+0.16} = \sqrt{0.25} = 0.5000,\\
\cos\theta &= 0.26/(0.5385\times0.5000) = 0.26/0.26926 = 0.9656, \qquad \theta = 15.07^\circ .
\end{align*}
\emph{After}, in $S_1$'s frame, $\vect{p}'_2=(0.5,1.0)$ and $\vect{p}'_3=(1.0,0.5)$:
\begin{align*}
\vect{p}'_2\cdot\vect{p}'_3 &= 0.5\times1.0+1.0\times0.5 = 1.0,\\
\norm{\vect{p}'_2} &= \norm{\vect{p}'_3} = \sqrt{0.25+1.0} = \sqrt{1.25} = 1.1180,\\
\cos\theta' &= 1.0/(1.1180\times1.1180) = 1.0/1.25 = 0.8000, \qquad \theta' = 36.87^\circ .
\end{align*}
Two spectra $15.07^\circ$ apart in reflectance are $36.87^\circ$ apart after the map\tagC. And $\vect{p}_1$, the darkest pixel
in both bands, became $(0,0)$: a vector of length zero, for which no angle is defined at all.
\end{worked}

Why, in general? Take a per-band map with no offset, $x'_b=s_bx_b$. Then
\begin{align*}
\cos\theta' &= \frac{\sum_b (s_bu_b)(s_bv_b)}{\sqrt{\sum_b s_b^2u_b^2}\,\sqrt{\sum_b s_b^2v_b^2}}
 && \text{substitute }u'_b=s_bu_b,\ v'_b=s_bv_b\\
&= \frac{\sum_b s_b^2u_bv_b}{\sqrt{\sum_b s_b^2u_b^2}\,\sqrt{\sum_b s_b^2v_b^2}} .
 && \text{collect } s_b\cdot s_b
\end{align*}
If every $s_b$ equals one number $s$, then $s^2$ comes out of each sum, $s^2/\sqrt{s^2}\sqrt{s^2}=1$, and
$\cos\theta'=\cos\theta$. If the $s_b$ differ, each band's product is weighted by its own $s_b^2$, and the angle changes.
With $s=(1/0.20,\,1/0.10)$, the pair above goes from $15.07^\circ$ to $9.25^\circ$\tagC. An offset $c_b$ changes it again,
because $\vect{u}+\vect{c}$ is no longer a multiple of anything related to $\vect{u}$. Min-max has both unequal scales
and offsets.

\begin{keypoint}
The fixed z-score changes angles too: it is also a per-band affine map. What it does not do is change them
\emph{differently in each scene}. A spectrum's z-scores are a property of the spectrum alone, identical for an image
pixel in any scene and for a library entry; that, not angle preservation, is what matching needs.
\end{keypoint}

%======================================================================
\section{A convolution over band index}\label{sec:hslnorm:conv}

The second obstacle sits in the tokenizer. Recall the convolutions of \cref{tab:hslwhat:trace}, step 1.

\begin{definition}[1-D convolution over band index; kernel; span]
A 1-D convolution with \emph{kernel} size $k$ makes band $i$'s output from bands $i-h,\dots,i+h$ with $h=(k-1)/2$,
using the same $k$ weights at every $i$. It knows the bands' \emph{positions in the list}, not their wavelengths. The
\emph{span} of the window at $i$ is $\lambda_{i+h}-\lambda_{i-h}$, the wavelength range its $k$ bands cover.
\end{definition}

The same $k$ weights are used everywhere, so the convolution assumes that ``two places along'' means the same
wavelength distance everywhere. On a list with roughly even spacing that is nearly true. On a list with a
water-absorption window removed it is false at the edges of the window.

\begin{worked}[label={we:hslnorm:span}]{A 5-band window on PRISMA's 149-band list}
PRISMA's 149-band list (the one every HyperSL experiment in this part used) drops 1325--1495\,nm. Its bands just below
and above that gap are, in nm\tagC:
\begin{align*}
\text{below: }&\dots,\ 1250.98,\ 1262.53,\ 1273.50,\ 1284.49,\ 1295.42,\ 1306.22,\ 1317.26,\\
\text{above: }&1502.02,\ 1512.63,\ 1523.22,\ \dots
\end{align*}
\emph{Window centred on 1273.50}, five bands 1250.98 to 1295.42:
\begin{align*}
\text{span} = 1295.42-1250.98 = 44.44\text{ nm}.
\end{align*}
\emph{Window centred on 1317.26}, five bands 1295.42, 1306.22, 1317.26, 1502.02, 1512.63. Its four steps are
\begin{align*}
1306.22-1295.42 &= 10.80,\\
1317.26-1306.22 &= 11.04,\\
1502.02-1317.26 &= 184.76,\\
1512.63-1502.02 &= 10.61,
\end{align*}
and its span is $1512.63-1295.42=217.21$\,nm\tagC. The kernel's third and fourth weights, which elsewhere connect two
bands 10\,nm apart, here connect 1317\,nm to 1502\,nm, across the 1.4\,$\mu$m water band.
\end{worked}

\Cref{tab:hslnorm:spans} gives the spans for the three kernels on four band lists, and \cref{fig:hslnorm:span} draws
the 5-band span along PRISMA's 149 bands.

\begin{table}[htbp]
\centering\small
\begin{tabular}{@{}p{50mm}rrrr@{}}
\toprule
band list & bands & $k=5$: median / max & $k=9$: max & $k=13$: median / max \\
\midrule
PRISMA, all 234 good bands (water kept) & 234 & 36.4 / 45.0 & 88.7 & 108.9 / 132.7 \\
PRISMA, 149 (water and noisy ranges cut) & 149 & 39.7 / 280.4 & 319.4 & 122.5 / 358.9 \\
EnMAP via SpectralEarth (water cut) & 202 & 31.0 / 291.6 & 335.0 & 93.1 / 379.2 \\
product 165-band grid & 165 & 40.0 / 200.0 & 240.0 & 120.0 / 280.0 \\
\bottomrule
\end{tabular}
\caption{How many nanometres (nm) a kernel over band index covers, for windows that fit inside the list\tagC. With the
water bands kept (first row; HyperSL trained on EnMAP this way) the spans are even: the design is consistent with its
own data. Every band list this product uses removes the water windows, and the largest span grows six- to sevenfold.}
\label{tab:hslnorm:spans}
\end{table}

\begin{figure}[htbp]
\centering
\begin{tikzpicture}
\begin{axis}[width=0.92\linewidth, height=55mm, xlabel={centre wavelength (nm)}, ylabel={span (nm)},
  ymin=0, ymax=380, xmin=450, xmax=2400, unbounded coords=jump, legend style={at={(0.02,0.98)},anchor=north west,font=\small}]
\addplot[spec, thick, mark=*, mark size=0.6pt] table[x=wl, y=span5] {figdata/hsl_kernel_span_prisma149.dat};
\addplot[loss, thick, mark=*, mark size=0.6pt] table[x=wl, y=span13] {figdata/hsl_kernel_span_prisma149.dat};
\legend{kernel 5, kernel 13}
\end{axis}
\end{tikzpicture}
\caption[The span of HyperSL's kernels along PRISMA's 149-band list]{The span of HyperSL's kernels along PRISMA's
149-band list. Away from the gaps a 5-band window covers about
40\,nm; next to the three removed ranges (1120--1186, 1317--1502 and 1785--2036\,nm) it jumps to 100--280\,nm. Data: \texttt{figdata/\allowbreak hsl\_kernel\_span\_prisma149.dat}, written by
\texttt{tools/\allowbreak compute\_hypersl.py}.}
\label{fig:hslnorm:span}
\end{figure}

\begin{sourcenote}[: the authors saw the effect and drew it]
The paper's Fig.~4 caption notes that ``the nonlinear segments in the figure are due to the removal of bad bands from
the ENMap data'' \citep{kong2025hypersl}. It refers to the position code, which knows wavelength and handles a gap
correctly. The convolutions in front of it do not; the code (\code{nn.Conv1d(1, 256, k, stride=1, padding='same')})
has no wavelength input.
\end{sourcenote}

Our own stem answers this differently: it groups bands into 40\,nm wavelength \emph{bins} (\cref{ch:specdecoder}), so a
band that is absent is simply absent, and no weight ever connects 1317\,nm to 1502\,nm as if they were neighbours.

%======================================================================
\section{What happened when it was tried}\label{sec:hslnorm:measured}

\subsection{The four-scene matching run (2026-09-23)}

The experiment \file{research/experiments/hypersl-splib-match/} (commits \code{f554fe4} to \code{4bbf91d}, no longer in
the working tree) matched every ground-truth anomaly pixel of the four labelled PRISMA scenes against two libraries:
ECOSTRESS (1,947 entries, 45 of them man-made) and USGS splib07 (518 entries)\tbrk\tagMr{the run's
\texttt{summary.json}, in \texttt{C:/Users/FITT/Documents/hypersl\_match\_data/results/}, 2026-09-23 22:02}. Three methods,
each ranked by cosine similarity:
\begin{itemize}
\item \textbf{HyperSL}: the embedding of pixel and library entry, each scaled \emph{per curve} (its own minimum and maximum
over bands), mask 0. Per-curve scaling was chosen because the scene frame clipped 41\% of the library and because a
spectrum cannot find itself across two different scalings (the control \code{frame_control.py});
\item \textbf{HyperSL raw}: the embedding of unscaled reflectance;
\item \textbf{SAM}: the cosine of the reflectances themselves, the classical reference.
\end{itemize}

\paragraph{Did anomalies match man-made entries?} The measure is the share of anomaly pixels whose best ECOSTRESS match is a
man-made entry, beside the same share for brightness-matched background pixels. A useful matcher should put the first well
above the second.

\begin{table}[htbp]
\centering\small
\begin{tabular}{@{}lrrrrrr@{}}
\toprule
& \multicolumn{2}{c}{HyperSL, per-curve} & \multicolumn{2}{c}{HyperSL, raw} & \multicolumn{2}{c}{SAM} \\
\cmidrule(lr){2-3}\cmidrule(lr){4-5}\cmidrule(lr){6-7}
scene (anomaly px) & anomaly & background & anomaly & background & anomaly & background \\
\midrule
set1 (638) & 7.2\% & 2.7\% & 0.16\% & 0.0\% & 9.6\% & 1.9\% \\
set2 (992) & 0.40\% & 0.0\% & 0.30\% & 0.0\% & 0.81\% & 0.20\% \\
set3 (778) & 0.26\% & 0.51\% & 0.39\% & 0.13\% & 2.7\% & 2.1\% \\
set4 (1,165) & 8.3\% & 33.0\% & 0.0\% & 0.0\% & 6.4\% & 24.3\% \\
\bottomrule
\end{tabular}
\caption{Share of pixels whose top ECOSTRESS match is man-made\tagMr{\texttt{summary.json}, \texttt{m1}}. HyperSL is below
SAM on three of four scenes, and on set4 both call a third or a quarter of ordinary background man-made.}
\label{tab:hslnorm:m1}
\end{table}

\paragraph{Where did the matches go?} The run also recorded the single library entry that took the most anomaly pixels
as its top match. With HyperSL per-curve, \emph{Salvia leucophylla 3}, a sage leaf, was the best match for 59.6\% of set2's
992 anomaly pixels, and \emph{Dolomitic Marble} for 31.2\% of set4's. Against USGS, a buckbrush spectrum took 52.5\% of
set2's. SAM's largest attractor on ECOSTRESS took 19.9\% (set2, \emph{Misc Grass Litter})\tagMr{\texttt{summary.json},
\texttt{top1\_share}}. An embedding in which one leaf is the nearest neighbour of most of a scene's anomalies is
describing the scene, not the anomalies.

\paragraph{Does the embedding at least organise the library?} For each library entry, take its 5 nearest other entries
and count the share in its own class:

\begin{center}\small
\begin{tabular}{@{}lrrrr@{}}
\toprule
& \multicolumn{2}{c}{ECOSTRESS} & \multicolumn{2}{c}{USGS} \\
\cmidrule(lr){2-3}\cmidrule(lr){4-5}
space & all classes & man-made & all classes & man-made \\
\midrule
HyperSL, per-curve & 0.965 & 0.124 & 0.827 & 0.373 \\
HyperSL, raw & 0.962 & 0.196 & 0.841 & 0.454 \\
SAM (reflectance) & 0.949 & 0.093 & 0.795 & 0.278 \\
\bottomrule
\end{tabular}
\end{center}
To be fair to HyperSL: its space groups library classes a little \emph{better} than raw angle, by 1.6 to 4.6 points
overall and 3 to 18 points for man-made entries\tagMr{\texttt{summary.json}, \texttt{purity}}. But look at the man-made
column on ECOSTRESS: 12\% means that seven out of eight of a man-made entry's nearest neighbours are natural. Vegetation,
minerals and soils dominate the library and the embedding; man-made materials are not a region of it.

\begin{sourcenote}[: why the USGS rows of the man-made table are all zero]
The run's man-made test is \code{e.chapter == "manmade"}. USGS splib07 calls its man-made chapter ``artificial'', so the
test is false for every USGS entry and every USGS share is 0 by construction (the library summary even reports
\code{"manmade": 0}). The purity measure uses \code{np.isin(cls, ["manmade", "artificial"])} and counts them. So read the
USGS purity, and ignore the USGS man-made shares; \cref{tab:hslnorm:m1} shows ECOSTRESS only for this reason.
\end{sourcenote}

\subsection{Band sampling moves the embedding as much as material does}

One more measurement shaped the decision, and it is the only one in this part whose raw output was not kept. On
2026-09-23 an ECOSTRESS entry fed at its own laboratory sampling found its \emph{own} PRISMA-resampled copy as nearest
neighbour only 28\% of the time\tbrk\tagM{docs/proposals/spectral-matching-foundation-model.md}. The proposal draws the
conclusion: ``that is the property a matching model must not have, and the reason not to start from HyperSL's
checkpoint''. The cloud-mask experiment of \cref{ch:hslcloud} took the lesson seriously and resampled every training
spectrum onto PRISMA's exact bands for that reason.

\begin{sourcenote}[: the 28\% has no output file]
The figure is quoted in the proposal (line 63) and in the team's memory note, both citing the 2026-09-23 run. The
experiment's committed scripts do not save it, and the run's \texttt{summary.json} does not contain it. Treat it as
recorded by the people who ran it, not as re-checkable. The related control in \code{frame_control.py} (can a pixel
find itself across two scalings?) printed its result to the screen and did not save it either.
\end{sourcenote}

%======================================================================
\section{What we do instead}\label{sec:hslnorm:ours}

\begin{center}\small
\begin{tabular}{@{}>{\raggedright\arraybackslash}p{34mm}>{\raggedright\arraybackslash}p{50mm}>{\raggedright\arraybackslash}p{50mm}@{}}
\toprule
requirement & HyperSL & this product \\
\midrule
one input per spectrum, in every scene & min-max of the scene: input depends on the scene's extremes &
\code{WavelengthStats}: one mean and spread curve over wavelength, fixed in the checkpoint \\
a library spectrum has an input & none; must borrow a frame (clips 41\%) or change the scaling &
the same curve read at the library's wavelengths, the same stem (\cref{ch:namer}) \\
brightness survives & removed per band, per scene & kept; the namer and ACE see reflectance \\
gaps in the band list & convolution over index joins 1317 and 1502\,nm & 40\,nm wavelength bins; an absent band is
absent \\
\bottomrule
\end{tabular}
\end{center}

\begin{drift}
The four-scene matching numbers belong to the 2026-09-23 run (ECOSTRESS 1,947 entries, USGS 518, PRISMA 149-band axis),
whose outputs live outside the repository. The span table depends on the band lists: PRISMA's from Set-1's HE5, the
SpectralEarth table, and the product grid as of 2026-10-31.
\end{drift}

%======================================================================
\begin{confusions}
\confusion{Min-max only changes brightness, and the spectral angle ignores brightness.}{It changes each band by its own
scale and offset; the angle changes (15.07° became 36.87° in \cref{we:hslnorm:angle}).}
\confusion{Then our z-score is just as bad.}{It is also affine, so it changes angles too, but identically in every scene and
for every library entry. Scene-dependence, not affinity, is the problem.}
\confusion{Scaling each spectrum by itself fixes it.}{It treats pixel and library alike, which is why the matching run used
it, but it is off the training distribution (the loss roughly doubles) and it erases brightness entirely.}
\confusion{The position code tells the tokenizer about the gap.}{The code is added after the convolutions; the convolutions
have already mixed 1317 and 1502\,nm.}
\confusion{HyperSL's embedding is worse than SAM at grouping the library.}{It is slightly better (0.965 against 0.949). It is
worse at the job that matters: anomalies matched to man-made entries.}
\confusion{The USGS man-made rate of 0\% means HyperSL never matched USGS man-made entries.}{It is 0 by construction: the
label test looked for ``manmade'' and USGS calls the chapter ``artificial''.}
\end{confusions}

\begin{checkyourself}
\item\label{ex:hslnorm:1} Scene $S_3$ has pixels $(0.15,0.25)$, $(0.20,0.50)$, $(0.25,0.55)$. Compute the min-max input of
  $(0.20,0.50)$ and compare with \cref{we:hslnorm:s1,we:hslnorm:s2}.
\item\label{ex:hslnorm:2} In $S_1$ plus the bright pixel $\vect{a}$, what is $\vect{p}_1$'s input? Why did it not move?
\item\label{ex:hslnorm:3} Compute the angle between $\vect{p}_2=(0.20,0.50)$ and $\vect{p}_3=(0.30,0.40)$ after the scaling
  $s=(1/0.20,\,1/0.10)$ with no offset. (The script gives 9.25°.)
\item\label{ex:hslnorm:4} Show that with $s_1=s_2=5$ and no offset the angle between any two 2-band spectra is unchanged.
\item\label{ex:hslnorm:5} On the product's 165-band grid the bands are 10\,nm apart except at the gaps; the 1790--1960\,nm gap
  removes the bands between. Find the span of the 5-band window whose centre is 1780\,nm, and the 13-band window centred
  at 1740\,nm.
\item\label{ex:hslnorm:6} A library spectrum is $(0.70,0.20)$. Write its fixed z-score with mean $(0.25,0.35)$ and spread
  $(0.10,0.12)$. Does the answer depend on which scene you are matching against?
\end{checkyourself}

\nextchapter{ch:hsldetect}
